\scanpage{3}
\mspage{\unsure{1}}
% archive foliation: 2
1. Ens. à quatre éléments.

Un ens. à quatre éléments $I$ est muni canoniquement d'une structure
\deleted{\unsure{d'espace vectoriel de}}\inserted{d'espace affine de}
\deleted{\unsure{de}} dim 2 sur $\mathbb{F}_2$, i.e. d'un sous-groupe
$V\subset \unsure{\mathcal{G}_I|\mathfrak{S}_I}$, isomorphe (non can.!) à $\mathbb{F}_2^2$,
tel que l'action de $V$ sur \deleted{\unsure{$\mathbb{F}$}} $I$ en fasse un torseur.
Le sous-groupe $V\subset \unsure{\mathcal{G}_I|\mathfrak{S}_I}$ est le seul sous-groupe
distingué abélien $\ne\{0\}$, \deleted{\unsure{c'est aussi}} il est son propre
\deleted{normalisateur}\inserted{centralisateur}, i.e.
$\unsure{\mathcal{G}_I|\mathfrak{S}_I}/V$ opère fidèlement sur $V$, en fait on a
\[
 \unsure{\mathcal{G}_I|\mathfrak{S}_I}/V
 \xrightarrow{\sim}\operatorname{Aut}_{\mathrm{gr}}(V)
 =\operatorname{Aut}_{\mathbb{F}_2\text{-vect}}(V)
 \xrightarrow{\sim}\deleted{\unsure{\operatorname{Aut}}}\unsure{\mathcal{G}_J|\mathfrak{S}_J}
\]
où $J=V^*$, la suite ex.te
\diagram{\begin{tikzcd}
1 \arrow[r] & V \arrow[r,dash] & \unsure{\mathcal{G}_I|\mathfrak{S}_I} \arrow[r] & \unsure{\mathcal{G}_J|\mathfrak{S}_J} \arrow[r] & 1
\end{tikzcd}}
étant \deleted{\unsure{isomorphe}}\inserted{identique} à la suite
\diagram{\begin{tikzcd}
1 \arrow[r,dash] & \underset{\text{translations}}{V} \arrow[r] &
\operatorname{Aut}_{\mathrm{aff}}(I) \arrow[r] & \operatorname{Aut}_{\mathbb{F}_2}(V) \arrow[r,dash] & 1 .
\end{tikzcd}}

Les \inserted{trois} éléments de $J=V^*$ correspondent aux trois
partitions de type $(2,2)$ de $I$ \inserted{ou aux trois involutions sans pt fixe}
(on peut l'interpréter \inserted{ainsi} car les éléments de $J$ correspondent
aux droites \inserted{vectorielles} de $V^*$, qui définissent
\deleted{\unsure{les}} donc des \deleted{\unsure{3}} relations d'équivalence
dans $I$, vers des espaces affines quotients de dim 1) --- les classes
intervenant dans l'ensemble de ces \inserted{3} relations d'équivalence
correspondent aux 6 \deleted{\unsure{paires}}\inserted{\unsure{parties}}
à deux éléments de $I$, ou encore aux 6 droites affines dans l'espace affine $I$,
l'application
\diagram{\begin{tikzcd}
\mathcal{P}_2(I) \arrow[r] & \mathcal{P}_{2,2}(I) \\
A \arrow[r,mapsto] & \{A,I\setminus A\}
\end{tikzcd}}
